\section{分次向量空间, 链复形与微分}\label{sec:graded-chain-complexes}

\emph{分次向量空间}是表示为以整数为指标的子空间 $V_k$ 的直和的向量空间 $V$:
\[
V=\bigoplus_{k=-\infty}^{\infty}V_k.
\]

例如, $n$ 个变量的多项式函数空间 $\mathcal P(\mathbb R^n)$ 是分次的: 当 $k\geq0$ 时, 第 $k$ 个直和项是次数为 $k$ 的齐次多项式空间, 否则为 $0$. 若分次向量空间像这个例子一样, 对 $k<0$ 有 $V_k=0$, 则称其为\emph{非负的}. 与全体多项式空间不同, 我们遇到的大多数分次空间都是非负的, 并且还是\emph{有限的}, 即只有有限个直和项非零. 例如, 次数至多为 $r$ 的多项式空间 $\mathcal P_r(\mathbb R^n)$ 是分次的, 并且只有 $r+1$ 个非零直和项.

分次向量空间之间的线性映射 $f:V\to W$, 若对 $k$ 有 $f(V_k)\subset W_{k+p}$, 则称为\emph{次数为 $p$ 的分次映射}. 例如, Laplace 算子在多项式上是次数为 $-2$ 的分次线性映射, 而乘以某个固定的次数为 $p\geq0$ 的齐次多项式, 则是次数为 $p$ 的分次映射.

\emph{链复形}就是配备了一个次数为 $-1$ 且满足 $\partial^2=0$ 的分次线性映射 $\partial:V\to V$ 的分次向量空间 $V$. 算子 $\partial$, 或等价地, 算子序列 $\partial_k=\partial|_{V_k}$, 称为链复形的\emph{微分}, 有时也称为\emph{边界算子}. 换言之, 链复形由一列向量空间 $V_k$ 和线性映射 $\partial_k:V_k\to V_{k-1}$, $k\in\mathbb Z$ 组成, 并满足 $\partial_k\circ\partial_{k+1}=0$. 我们常将链复形写成如下形式:
\[
\cdots\to V_{k+1}\xrightarrow{\partial_{k+1}}V_k\xrightarrow{\partial_k}V_{k-1}\to\cdots,
\]

但在需要简略表示时, 则只写成 $(V,\partial)$, 甚至只写成 $V$. $V_k$ 的元素有时称为 $k$-链.

分次向量空间的\emph{分次子空间}, 自然是指可以分解为子空间 $S_k\subset V_k$ 的直和的子空间 $S$. 若 $S$ 是复形 $(V,\partial)$ 的分次子空间, 且 $\partial S_k\subset S_{k-1}$, 则 $(S_k,\partial_k|_{S_k})$ 本身也是一个复形, 称为 $(V,\partial)$ 的\emph{子复形}.
